\section{Introduction}
\label{sec:intro}

Modern machine learning research is characterized by rapid, combinatorial
growth: advances in reasoning~\cite{wei2022chain, yao2023tree}, tool
use~\cite{schick2023toolformer, qin2023toolllm}, and multi-modal
integration~\cite{shen2023hugginggpt, lu2023chameleon} each introduce method
components that practitioners must understand, compare, and selectively combine
to push the state of the art. A single competitive submission to a top venue
now routinely assembles half a dozen individually published techniques---chain-of-thought prompting layered with tree search, tool-augmented verification,
self-consistency filtering, and domain-specific adaptations---into an
integrated pipeline. As the space of available building blocks expands, the
cognitive burden of \emph{method composition} becomes a bottleneck: researchers
must simultaneously track technical compatibility, identify complementary
strengths, avoid known failure interactions, and spot the gaps where novel
contributions are needed. This complexity creates a compelling opportunity for
computational support in the method design process itself.

\paragraph{Limitations of existing approaches.}
Recent work has explored two broad strategies for automating aspects of
research. On one end, \emph{end-to-end idea generation} systems such as The AI
Scientist~\cite{lu2024aiscientist} and AI Co-Scientist~\cite{gottifredi2024aicoscientist} prompt large language models to
produce full research proposals or even complete papers. While impressive in
scope, these systems treat idea generation as monolithic text completion,
lacking explicit representations of method structure, compatibility
constraints, or compositional reasoning. As a result, they frequently produce
ideas that are either superficial recombinations of known concepts or
internally inconsistent designs that overlook critical failure
modes~\cite{huang2024mlrbench}. On the other end, \emph{knowledge retrieval}
approaches---including semantic search engines, citation graphs, and concept
extraction pipelines~\cite{baek2024researchagent, li2024chainofideas}---excel
at surfacing relevant prior work but stop short of composing retrieved
knowledge into actionable method designs. They organize \emph{what is known}
without reasoning about \emph{what could be built}. Neither paradigm provides
the structured, constraint-aware composition that effective method design
demands.

\paragraph{Key insight: methods as composable skill primitives.}
Our work is motivated by a simple but powerful observation: the methodological
contributions in research papers can be decomposed into \emph{research
skills}---modular primitives with well-defined functional interfaces. Each
skill specifies a goal (what problem it addresses), a mechanism (how it works),
assumptions (when it applies), typed inputs and outputs, empirical performance
gains, known failure modes, and explicit compatibility and conflict
declarations with other skills. This structured decomposition is analogous to
the shift from monolithic programs to modular software with typed APIs: once
methods are represented as interface-bearing components, systematic composition
becomes tractable. Skills can be organized into a graph whose edges encode
functional relationships---prerequisite dependencies, enhancement
opportunities, substitution alternatives, conflict constraints---enabling
compositional search that respects the rich structure of methodological
knowledge.

This perspective draws conceptual inspiration from compositional frameworks in
adjacent domains. In embodied AI, Voyager~\cite{wang2023voyager} builds a
library of reusable behavioral skills for open-ended exploration. In tool-use
research, systems such as ToolLLM~\cite{qin2023toolllm} and
HuggingGPT~\cite{shen2023hugginggpt} compose API-bearing tools into
multi-step pipelines. ResearchOS extends this compositional philosophy from
executable tools and game behaviors to \emph{research methodology itself},
treating the intellectual building blocks of scientific method design as
first-class composable objects.

\paragraph{Our approach: ResearchOS.}
We introduce ResearchOS, a framework that operationalizes skill-based method
composition through a five-stage pipeline (Figure~\ref{fig:pipeline}):
%
\begin{enumerate}[nosep,leftmargin=*]
  \item \textbf{Skill Extraction.} Given a corpus of research papers, an LLM-based extraction module identifies and structures each methodological
  contribution as a research skill with a standardized interface schema.
  \item \textbf{Skill Graph Construction.} Extracted skills are organized into
  a \emph{Research Skill Graph} with six typed relation edges---prerequisite,
  enhances, substitutes, conflicts, composes, and refines---that encode
  functional dependencies and opportunities among skills.
  \item \textbf{Constraint-Aware Composition.} Given a research goal, a
  composition engine performs goal decomposition, retrieves candidate skills
  from the graph, assembles them into a directed acyclic graph (DAG)
  respecting compatibility constraints, performs gap analysis to identify
  missing capabilities, and proposes novel bridging skills to fill those gaps.
  An iterative self-critique loop refines the design for internal consistency.
  \item \textbf{Multi-Level Output Generation.} The composed design is
  rendered at three levels of specificity: a \emph{Method Design Card}
  (structured summary), a \emph{Research Proposal} (narrative with
  motivation, hypotheses, and expected results), and an \emph{Executable
  Experiment Plan} (concrete implementation steps, baselines, and evaluation
  protocol).
  \item \textbf{Validation.} Generated designs are assessed through a
  multi-level evaluation protocol combining automated quality metrics,
  pairwise Elo ranking by LLM judges, and future-grounded novelty analysis.
\end{enumerate}

\paragraph{Contributions.}
This paper makes four principal contributions:
\begin{enumerate}[nosep,leftmargin=*]
  \item We introduce the \textbf{research skill schema}, a structured
  representation of methodological primitives with typed functional interfaces
  (goal, mechanism, assumptions, inputs/outputs, empirical gains, failure
  modes, and compatibility/conflict declarations) that enables systematic
  decomposition of published methods into composable units
  (Section~\ref{sec:skill_schema}).

  \item We propose the \textbf{Research Skill Graph}, a heterogeneous
  knowledge graph with six relation types (prerequisite, enhances, substitutes,
  conflicts, composes, refines) that captures the functional landscape of a
  research domain and supports structured retrieval and constraint
  checking (Section~\ref{sec:skill_graph}).

  \item We develop a \textbf{constraint-aware composition engine} that
  assembles skills into novel method designs through goal decomposition,
  DAG construction with constraint satisfaction, and a gap-bridging mechanism
  that identifies and fills missing capabilities with novel skill proposals,
  followed by iterative self-critique
  (Section~\ref{sec:composer}).

  \item We design a \textbf{multi-level evaluation protocol} combining
  pairwise Elo ranking, future-grounded novelty assessment via temporal
  corpus splits, small-scale execution validation, and systematic judge
  reliability analysis, providing rigorous and multi-faceted assessment of
  generated research designs (Section~\ref{sec:experiments}).
\end{enumerate}

\paragraph{Results preview.}
We instantiate ResearchOS on the domain of LLM reasoning, extracting over 200
skills from more than 100 recent papers. In head-to-head comparisons,
ResearchOS designs achieve significantly higher Elo ratings than those produced
by direct prompting, retrieval-augmented generation (RAG), and flat skill list
baselines. Future-grounded evaluation using temporal corpus splits
demonstrates that ResearchOS-generated designs anticipate methodological
directions that appear in subsequent real publications, confirming genuine
novelty rather than mere recombination~\cite{du2024hindisight}. Ablation
studies reveal that both the graph structure (versus unstructured skill lists)
and the gap-bridging mechanism (versus composition without gap filling) are
critical to design quality, with their removal causing substantial drops in
novelty and feasibility scores. These results establish that structured,
constraint-aware skill composition offers a principled and effective paradigm
for automated research method design.
